\documentclass{article}
\usepackage[UTF8]{ctex}
\usepackage{amsmath,amssymb,amsfonts}
\usepackage{enumitem}
\usepackage{geometry}
\geometry{a4paper, margin=1in}
% === 错题本排版优化（不参与题目标识） ===
\setlength{\parindent}{0pt}
\setlength{\parskip}{0.35em plus 0.10em minus 0.08em}
\setlength{\emergencystretch}{3em}
\setlist{itemsep=0.25em, topsep=0.40em, parsep=0pt}
\allowdisplaybreaks[2]
% === 错题本排版优化结束 ===
\begin{document}
\section*{中国矿业大学2025年数学分析真题}
\section*{23．中国矿业大学2025年数学分析试卷}

\subsection{题目1}
一．（15 分）已知函数 \(f(x)\) 是区间 \([a, b]\) 上的单调递增函数，且 \(f(a)>a, f(b)<b\) ．证明：存在一个数\(\xi \in(a, b)\) ，使得 \(f(\xi)=\xi\) ．

\subsection{题目2}
二．（15 分）已知函数 \(f(x)=\frac{2 x}{1+x^{2}}\) ，求 \(f(x)\) 的单调区间，极值点，凹凸区间，渐近线，并作出大致草图．

\subsection{题目3}
三．（15 分）已知 \(f(x)\) 是 \([a, b]\) 上的连续函数，\(x_{1}<x_{2}<\cdots<x_{n}\) 是 \([a, b]\) 上的点．

\begin{itemize}
  \item[（1）] （5 分）叙述闭区间上的连续函数介值定理．
  \item[（2）] （10 分）证明：存在 \(\xi \in\left[x_{1}, x_{n}\right]\) ，使得 \(f(\xi)=\sum_{j=1}^{n} \frac{f\left(x_{j}\right)}{n}\) ．
\end{itemize}

\subsection{题目4}
四．（15 分）求下列不定积分：

\begin{itemize}
  \item[（1）] （5 分） \(\int \frac{1}{e^{x}-e^{-x}} \mathrm{~d} x\) ．
  \item[（2）] （10 分） \(\int \frac{1}{x+\sqrt{x^{2}-x+1}} \mathrm{~d} x\) ．
\end{itemize}

\subsection{题目5}
五．（15 分）已知 \(f(x)\) 在 \([a, b]\) 上可积．

\begin{itemize}
  \item[（1）] （5 分）若 \(f(x)\) 是 \([a, b]\) 上的连续非负函数，且存在 \(x_{0} \in[a, b]\) ，使得 \(\lim _{x \rightarrow x_{0}} f(x)=A>0\) ．证明：
\[
\int_{a}^{b} f(x) \mathrm{d} x>0 .
\]
  \item[（2）] （10 分）若 \(\int_{a}^{b} f(x) \mathrm{d} x>0\) ，证明：存在闭区间 \([\alpha, \beta] \subset[a, b]\) ，使得 \(f(x)\) 在 \([\alpha, \beta]\) 上满足 \(f(x)>0\) ．
\end{itemize}

\subsection{题目6}
六．（15 分）已知 \(a_{n}=\int_{0}^{\frac{\pi}{4}}(\tan x)^{n} \mathrm{~d} x\) ，其中 \(n=1,2, \cdots\) ．

\begin{itemize}
  \item[（1）] （5 分）求级数 \(\sum_{n=1}^{\infty} \frac{a_{n}+a_{n+2}}{n}\) 的和．
  \item[（2）] （10 分）证明：对任意常数 \(p>0\) ，数项级数 \(\sum_{n=1}^{\infty} \frac{a_{n}}{n^{p}}\) 都收敛。
\end{itemize}

\subsection{题目7}
七．（15 分）求重积分 \(I=\iint_{D} \frac{3 x}{y^{2}+x y^{3}} \mathrm{~d} x \mathrm{~d} y\) ，其中 \(D=\left\{(x, y) \in \mathbb{R}^{2} \mid 1 \leq x y \leq 3, x \leq y^{2} \leq 3 x\right\}\) ．

\subsection{题目8}
八．（15 分）已知 \(f(x, y)\) 在 \(P_{0}\left(x_{0}, y_{0}\right)\) 的某邻域 \(U\left(P_{0}\right)\) 上有二阶连续偏导数，且 \(f_{x}\left(x_{0}, y_{0}\right)=\) \(0, f_{x x}\left(x_{0}, y_{0}\right)>0\) 。证明：存在关于 \(y_{0}\) 的某邻域 \(V\left(y_{0}\right)\) ，对任意 \(y \in V\left(y_{0}\right), f(x, y)\) 作为关于 \(x\)的函数存在极小值 \(g(y)\) ，且 \(g^{\prime}\left(y_{0}\right)=f_{y}\left(x_{0}, y_{0}\right)\) 。

\subsection{题目9}
九。（15 分）已知 \(u(x, y)\) 在 \(\mathbb{R}^{2}\) 上有二阶连续偏导数，若 \(u_{x x}+u_{y y}=0\) 对任意 \((x, y) \in \mathbb{R}^{2}\) 都成立，则称 \(u(x, y)\) 为 \(\mathbb{R}^{2}\) 上的调和函数。证明：\(u(x, y)\) 为 \(\mathbb{R}^{2}\) 上的调和函数的充要条件是 \(\oint_{C} \frac{\partial u}{\partial n} \mathrm{~d} s=0\) ，其中\(C\) 为 \(\mathbb{R}^{2}\) 上任意光滑闭曲线， \(\mathbf{n}\) 为 \(u(x, y)\) 沿 \(C\) 的单位外法线方向。

\subsection{题目10}
十．（15 分）已知欧拉积分 \(\int_{0}^{+\infty} e^{-x^{2}} \mathrm{~d} x=\frac{\sqrt{\pi}}{2}\) ，求含参量积分 \(I(a)=\int_{0}^{+\infty} e^{-\left(x^{2}+\frac{a^{2}}{x^{2}}\right)} \mathrm{d} x\) 的值，其中\(a>0\) ．
\end{document}
